\bookchapter{Capitolo 5: Rendi il telefono più simile a te — Impostazioni di base}

\chapterintro{Qualche minuto di impostazioni e l'esperienza migliora enormemente. Sono come “accordare” il telefono, per adattarlo alle tue abitudini d'uso.}

\booksection{Suono e vibrazione}

Imposta una suoneria e un suono di notifica che ti piacciono. Usa con flessibilità la modalità silenziosa (in riunione) e la modalità non disturbare (per dormire): il mondo diventa subito più silenzioso.

\begin{itemize}
  \item \textbf{Suoneria}: scegline una che riesci a sentire e non imbarazzante.
  \item \textbf{Modalità vibrazione}: in silenzioso, avvisa con la vibrazione.
  \item \textbf{Modalità non disturbare}: puoi impostare eccezioni, ad esempio consentire le chiamate dai contatti.
  \item \textbf{Volumi separati}: volume multimediale, suoneria e sveglia si possono regolare separatamente.
\end{itemize}

\booksection{Display e luminosità}

\begin{itemize}
  \item \textbf{Luminosità automatica}: lascia che il telefono si regoli in base alla luce ambientale, senza pensieri.
  \item \textbf{Modalità riposo occhi / Night Shift}: filtra la luce blu, più confortevole la sera.
  \item \textbf{Modalità scura}: risparmia energia (soprattutto sugli schermi OLED) e di notte non abbaglia. Molte app la supportano.
  \item \textbf{Frequenza di aggiornamento dello schermo}: alcuni Android permettono di regolarla; più alta (es. 120 Hz) = scorrimento più fluido.
  \item \textbf{Blocco automatico}: imposta 15 o 30 secondi, sicuro e risparmia batteria.
  \item \textbf{Dimensione dei caratteri}: ti sembrano piccoli? Ingrandiscili.
\end{itemize}

\booksection{Wallpaper, temi e home}

Cambia il wallpaper con uno bello (statico o animato) e l'umore cambia. I telefoni Android di marca hanno spesso un tema store che permette di cambiare icone e stile in un colpo solo.

\textbf{Widget sulla home}: metti meteo, agenda e controllo musica sulla home, per arrivarci in un tocco. Lo supportano sia iOS che Android; tieni premuto su un'area vuota della home per aggiungerli.

\booksection{Impostazioni di efficienza: trasforma il telefono nel tuo “cervello esterno”}

Sono funzionalità “una volta impostate non torni più indietro”: ciascuna può migliorare sensibilmente l'efficienza quotidiana.

\booksubsection{Modalità Focus / Non disturbare: riprenditi l'attenzione rubata dalle notifiche}

\begin{itemize}
  \item \textbf{Idea centrale}: non "guardare meno il telefono", ma "concentrarsi quando serve, farsi sentire quando si può".
  \item \textbf{Suggerimenti d'uso}:\newline · Modalità lavoro: consenti solo contatti "speciale" e app indispensabili;\newline · Modalità sonno: attivazione automatica, consenti solo la sveglia;\newline · Allenamento / lettura / riunione: personalizza in base all'attività corrente.
  \item \textbf{Uso avanzato}: usa automazione + posizione/orario per cambiare modalità in automatico. Ad esempio, Comandi rapidi su iOS può impostare "all'arrivo in ufficio entra automaticamente in modalità lavoro".
\end{itemize}

\booksubsection{Gestione delle notifiche: basta “assilli” continui dai pallini rossi}

\begin{itemize}
  \item \textbf{Disattiva le notifiche marketing}: il 90\% delle push in realtà non ti serve.
  \item \textbf{Ordina per importanza}: metti in cima le comunicazioni interpersonali (telefono, SMS, IM); notizie, shopping e social possono essere "notifiche silenziose".
  \item \textbf{Raggruppa le notifiche}: piega insieme le notifiche dello stesso tipo, per non essere sommerso da decine di messaggi a ogni sblocco.
\end{itemize}

\booksubsection{Gesti rapidi: operazioni un passo avanti}

\begin{itemize}
  \item \textbf{Screenshot}: tasto laterale + volume giù (iOS full-screen) / tasto power + volume giù (Android generale).
  \item \textbf{Screenshot lungo / scorrevole}: dopo lo screenshot normale, tocca l'opzione "screenshot lungo".
  \item \textbf{Scansione QR con un tocco}: aggiungi l'icona "scansione" nel centro di controllo / impostazioni rapide.
  \item \textbf{Tocca sul retro (iPhone)}: in [Impostazioni] > [Accessibilità] > [Tocco] > [Tocca sul retro] puoi associare "screenshot" o "blocco schermo" ad azioni di doppio/triplo tocco.
  \item \textbf{Risveglio vocale}: Siri / Xiao Ai Tongxue / Xiao Yi / Xiao Bu: basta chiamare per conoscere il meteo o impostare la sveglia.
\end{itemize}

\booksubsection{Tastiera: l'arma di efficienza più sottovalutata}

\begin{itemize}
  \item \textbf{Scegli la tastiera giusta}:\newline · Sogou / Baidu: dizionario ampio, aggiornamento rapido di slang e parole di moda;\newline · iFlytek: il miglior riconoscimento vocale, perfetta per prendere appunti in riunione;\newline · Gboard (Google): minimalista, forte sincronia multipiattaforma, supporta digitazione scorrendo;\newline · Tastiera di iOS: la migliore per privacy, integrata in profondità nel sistema.
  \item \textbf{Interruttori indispensabili}: correzione intelligente, suggerimenti, sincronia del dizionario cloud, scorrimento verso l'alto per numeri e simboli.
  \item \textbf{Nota sulla privacy}: usa con cautela le funzioni di "suggerimento per intere frasi" o "raccomandazioni personalizzate", potrebbero caricare i contenuti digitati.
\end{itemize}

\booksubsection{Home e multitasking: rendi fluido il flusso di lavoro}

\begin{itemize}
  \item \textbf{Split-screen / Picture-in-Picture}: rispondi ai messaggi mentre guardi un video, prendi appunti mentre navighi.
  \item \textbf{Finestre piccole / fluttuanti}: WeChat / Douyin supportano la "modalità finestra piccola", richiudila dopo aver risposto.
  \item \textbf{App più usate in evidenza}: nel Dock (in basso nello schermo) metti solo le 4 app ad alta frequenza, per ridurre il tempo di ricerca.
\end{itemize}

\begin{tcbinfo}{Checklist "primo giorno di lavoro"}
Con un telefono nuovo o dopo un grande aggiornamento, dedica 10 minuti a queste impostazioni:\newline (1) Password di blocco + biometria (2) Raggruppamento notifiche + disattivazione push marketing (3) Dimensione caratteri / frequenza di aggiornamento / modalità scura\newline (4) Tastiera / browser / client mail predefiniti (5) Contatti di emergenza / cartella clinica (6) Impostazioni di backup

\end{tcbinfo}

\booksection{Accessibilità e inclusività: la forza della tecnologia per il bene}

Non farti intimidire dal nome “accessibilità”: queste funzioni sono per tutti. Inizialmente pensate per persone con disabilità, sono utili anche per gli altri — mani occupate con i pacchi, cuffie rotte in metropolitana, caratteri piccoli con gli occhiali da vista… possono salvare la situazione. In fondo, il design accessibile riflette la filosofia della\textbf{tecnologia per il bene}: permettere a chiunque di godere in modo paritario della vita digitale.

\booksubsection{Gli “occhi” degli ipovedenti — Guida completa agli screen reader}

Lo screen reader è lo strumento centrale per gli ipovedenti:\textbf{legge ad alta voce tutto il contenuto visibile a schermo}e l'utente opera tramite gesti e voce. Qui introduciamo l'uso di VoiceOver di iOS e TalkBack di Android.

\begin{tcbinfo}{VoiceOver di iOS}
\textbf{Percorso di attivazione}: [Impostazioni] > [Accessibilità] > [VoiceOver], attiva l'interruttore. Oppure triplo clic sul tasto laterale (o Home) per attivare/disattivare rapidamente.

\textbf{Gesti di base}:

\begin{enumerate}
  \item Scorri con un dito: esplora il contenuto elemento per elemento;
  \item Scorri con due dita: scorri la pagina intera su/giù o sinistra/destra;
  \item Doppio tocco con un dito: attiva l'elemento attivo (equivalente al "tap");
  \item Scorri con tre dita: cambia pagina;
  \item Triplo tocco con un dito: zoom o selezione (diverso dal triplo tap normale).
\end{enumerate}

\textbf{Personalizzazione}: [VoiceOver] > [Voce] permette di regolare la\textbf{velocità di lettura}(default circa 50\%, portarla a 70-80\% per un ascolto più fluido), il\textbf{tono}e di cambiare voce. iOS include “Voce avanzata” e “Controllore vocale”: la prima riduce respiri e pause, per un ascolto più naturale.

\end{tcbinfo}

\begin{tcbinfo}{TalkBack di Android}
\textbf{Percorso di attivazione}: [Impostazioni] > [Accessibilità] > [TalkBack], attiva l'interruttore. Oppure tieni premuti i tasti volume su+giù per 3 secondi (su alcuni brand).

\textbf{Gesti di base}:

\begin{enumerate}
  \item Scorri a sinistra/destra con un dito: cambia elemento attivo;
  \item Doppio tocco con un dito: attiva l'elemento;
  \item Scorri su/giù o sinistra/destra con due dita: scorri;
  \item Tap con tre dita: entra in modalità lettura (lettura continua);
  \item Tap con quattro dita: pausa/riprendi la lettura.
\end{enumerate}

\textbf{Personalizzazione}: [TalkBack] > [Impostazioni text-to-speech] permette di regolare\textbf{velocità}e\textbf{tono}, nonché di cambiare\textbf{motore vocale}(consigliato il motore Google Text-to-Speech per una lettura più naturale).

\end{tcbinfo}

\booksubsection{Da piccolo a grande, da sfocato a nitido — Adattamenti visivi}

\booksubsection{Gesti di ingrandimento}

\begin{itemize}
  \item \textbf{iOS}: [Impostazioni] > [Accessibilità] > [Zoom], una volta attivo\textbf{tre tocchi leggeri sullo schermo}ingrandiscono. Di default ingrandisce a schermo intero, è possibile scegliere zoom in finestra.\textbf{Pizzica con due dita per regolare il fattore di zoom (da 1,5x a 15x)}; doppio tocco con tre dita per bloccare lo zoom corrente.
  \item \textbf{Android}: [Impostazioni] > [Accessibilità] > [Ingrandimento], tocca con tre dita per ingrandire, di nuovo per uscire; oppure\textbf{tre tocchi con tre dita}per uno zoom temporaneo (massimo 8x).
\end{itemize}

\booksubsection{Testo ad alto contrasto / grassetto}

\begin{itemize}
  \item \textbf{iOS}: [Impostazioni] > [Accessibilità] > [Display e dimensione testo], attiva "Grassetto", "Aumenta contrasto", "Riduci trasparenza". "Inversione intelligente" mantiene il testo nero su sfondo chiaro, evitando testo bianco su bianco.
  \item \textbf{Android}: [Impostazioni] > [Display] > [Dimensione caratteri], regolabile su "Piccolo, Normale, Grande, Enorme". Huawei/Xiaomi offrono anche i toggle "Grassetto" e "Testo ad alto contrasto".
\end{itemize}

\booksubsection{Filtri colore (per daltonismo)}

\begin{itemize}
  \item \textbf{iOS}: [Impostazioni] > [Accessibilità] > [Display e dimensione testo] > [Filtri colore], offre\textbf{quattro modalità daltonismo rosso/verde/blu/giallo}, più "Scala di grigi" e "Inversione colori". Consigliato un test della vista (integrato) prima di scegliere il filtro.
  \item \textbf{Android}: [Impostazioni] > [Accessibilità] > [Correzione colore], include\textbf{tre modalità daltonismo}+ il toggle "Inversione colori".
\end{itemize}

\booksubsection{Descrizione delle immagini con AI: l'AI come “occhio”}

Quando lo screen reader incontra un'immagine, la pratica tradizionale è saltarla o leggere solo la parola “immagine”. Ma oggi l'AI può\textbf{descrivere il contenuto dell'immagine in una frase}, molto utile per riconoscere captcha, leggere grafici e menu.

\begin{tcbinfo}{Supporto alla descrizione AI delle immagini sui sistemi}
\begin{itemize}
  \item \textbf{iOS 15+}: in [Impostazioni] > [Accessibilità] > [VoiceOver] > [Descrizioni vocali] attiva "Descrizioni immagini". Il sistema usa un\textbf{modello AI locale}per generare automaticamente la descrizione testuale (richiede iPhone 12 o successivo).
  \item \textbf{Nuova app TalkBack}: le versioni più recenti di TalkBack integrano la funzione\textbf{Q\&A sull'immagine}: puoi chiedere "cosa c'è in questa immagine?" e l'AI risponde in tempo reale (non dipende strettamente dalla versione di Android, basta aggiornare TalkBack).
  \item \textbf{App di terze parti}: "Seeing AI" (iOS), "Lookout" (Android), "Xiao Yi Bang Bang" (Huawei) sono eccellenti assistenti visivi AI: riconoscono oggetti, leggono documenti, identificano espressioni facciali.
\end{itemize}

\end{tcbinfo}

\begin{tcbpractice}{Scenario d'uso}
Prova a chiedere di una foto di strada “quante persone ci sono in questa immagine? Cosa stanno facendo?”, oppure apri un'app di delivery e fai leggere dall'AI il nome dei piatti nel menu —\textbf{questa è la vera “tecnologia che rende la vita più paritaria”}.

\end{tcbpractice}

\booksubsection{Le “orecchie” dei non udenti — Guida ai sottotitoli in tempo reale}

I sottotitoli live (Live Captions) sono una funzione\textbf{di sistema di speech-to-text}, molto utili per i non udenti in ambienti rumorosi, in biblioteca o in contesti con lingue straniere.

\begin{tcbinfo}{Attivazione e ottimizzazione dei sottotitoli live}
\textbf{iOS}: [Impostazioni] > [Accessibilità] > [Sottotitoli live], attiva. Il sistema rileva tutti gli audio del dispositivo (chiamate, video, podcast, FaceTime, Memo vocali) e mostra i sottotoli in alto sullo schermo.

\textbf{Android 12+}: premi il\textbf{tasto volume}tocca l'icona “Sottotitoli live” nel popup; oppure [Impostazioni] > [Accessibilità] > [Sottotitoli live].

\textbf{Consigli per migliorare l'accuratezza}:

\begin{enumerate}
  \item usali preferibilmente in un\textbf{ambiente silenzioso}, lontano da condizionatori e rumori meccanici;
  \item l'accuratezza è massima quando chi parla\textbf{articola bene e mantiene un ritmo medio};
  \item le conversazioni a più voci riducono l'accuratezza, puoi usare la modalità "mostra solo una persona";
  \item supporta lingue principali come cinese mandarino, cantonese, inglese;\textbf{lingue minori hanno accuratezza più bassa};
  \item musica di fondo o voci interferenti rendono il riconoscimento discontinuo.
\end{enumerate}

\end{tcbinfo}

\booksubsection{Doppio canale in chiamata: speech-to-text + text-to-speech}

Quando l'utente\textbf{non sente e non può parlare}(afonia improvvisa, post-operatorio alla gola, balbuzie grave), il telefono può fare da “interprete”.

\begin{tcbinfo}{Trascrizione in tempo reale della chiamata (TTY / RTT)}
\textbf{iOS}: [Impostazioni] > [Accessibilità] > [RTT/TTY], una volta attivo, in chiamata compare una tastiera a schermo; la voce dell'interlocutore viene trascritta in testo in tempo reale e\textbf{quello che digiti viene letto a voce}all'interlocutore.

\textbf{Android}: [Impostazioni] > [Accessibilità] > [Trascrizione live] + [Avvisi sonori]. “Trascrizione live” mostra il testo di chi parla; “Avvisi sonori” converte in voce il tuo input quando non puoi parlare.

\textbf{Compatibilità}:

\begin{itemize}
  \item \textbf{RTT (Real-Time Text)}: entrambi gli interlocutori devono supportarlo, il testo compare mentre si digita, senza ritardo; già supportato dai principali operatori;
  \item \textbf{TTY (telescrivente)}: richiede un dispositivo TTY dedicato, sempre più sostituito da RTT;
  \item la\textbf{funzione voice-to-text}di WeChat/QQ è un'alternativa: tieni premuto sul messaggio vocale "Converti in testo".
\end{itemize}

\end{tcbinfo}

\booksubsection{Funzioni “piccole ma utili” per tutti}

\begin{itemize}
  \item \textbf{Tocca sul retro}(iOS 14+): [Impostazioni] > [Accessibilità] > [Tocco] > [Tocca sul retro], puoi associare a\textbf{doppio/triplo tocco}azioni come screenshot, blocco schermo, centro di controllo. Particolarmente comodo con le mani occupate dai pacchi.
  \item \textbf{AssistiveTouch / pallina sospesa}: iOS [Accessibilità] > [Tocco] > [AssistiveTouch]; Android [Accessibilità] > [Navigazione sospesa]. Un pulsante virtuale che sostituisce gesti o tasti fisici, utile quando i tasti non funzionano o si usa il telefono con una sola mano.
  \item \textbf{Controllo interruttori}: controlla il telefono tramite interruttori esterni, tap singolo o movimenti della testa,\textbf{pensato per utenti con disabilità motorie}. iOS [Accessibilità] > [Controllo interruttori]; Android [Accessibilità] > [Controllo interruttori].
  \item \textbf{Riconoscimento suoni}(iOS): [Accessibilità] > [Riconoscimento suoni], una volta attivo il telefono "ascolta" campanello, pianto di neonato, allarme antincendio, sirene e\textbf{alla rilevazione vibra + notifica a schermo}; molto utile per i non udenti e durante il sonno.
  \item \textbf{Avviso LED lampeggiante}: iOS [Accessibilità] > [Audio/Visivo] > [Lampeggio LED per avvisi]; Android [Impostazioni] > [Notifiche] > [Avviso flash]. A chiamate e notifiche il flash lampeggia,\textbf{percezione senza suono anche in ambienti rumorosi}.
  \item \textbf{Audio mono}: [Impostazioni] > [Accessibilità] > [Audio] > [Audio mono], unisce i canali sinistro e destro. Chi ha un orecchio con calo uditivo non perde metà delle informazioni sonore.
\end{itemize}

\begin{tcbinfo}{La filosofia di accessibilità di Apple}
Apple investe moltissimo nell'accessibilità e ogni anno aggiorna e aggiunge funzioni. Sul sito Apple c'è una pagina dedicata all'accessibilità, con spiegazioni dettagliate. Le funzioni di accessibilità di iOS sono in genere più complete e fluide rispetto ad Android, motivo per cui molti utenti con disabilità scelgono Apple.

\end{tcbinfo}

\begin{tcbpractice}{30 secondi per provare}
Trova un momento libero e fai queste prove:

\begin{enumerate}
  \item Attiva\textbf{VoiceOver / TalkBack}e prova a scorrere WeChat con tre dita, senza usare gli occhi;
  \item Attiva i\textbf{sottotitoli live}e parla al telefono per valutare l'accuratezza;
  \item \textbf{Tocca con tre dita}per ingrandire un testo al 200\% e vivere la quotidianità di un ipovedente;
  \item Disattiva tutto e pensa: "se avessi sempre bisogno di queste funzioni, come vorrei che fossero progettate".
\end{enumerate}

\end{tcbpractice}
